% Lecture example set: SC3000-L03
% Sources: CZ3005 Module 3 pp. 32--39; CZ3005 Module 4 pp. 4--7; discussion check
% Human verified: Yes

\clearpage
\exercisemeta
  {SC3000-L03-EX}
  {2026-08-31}
  {Module 3 pp.~32--39；Module 4 pp.~4--7；课堂讨论补充}
  {答案已根据讲义与讨论核对}
  {已确认（2026-08-31）}

\exercisequestion{SC3000-L03-E01}{Minimax Tree 的逐层回传}

\textbf{对应知识点：}
\begin{itemize}
  \item \relatedtopic{SC3000-L03-T01}{Game Formulation 与 Minimax}
\end{itemize}

\subsubsection*{题目}

在\cref{fig:sc3000-l03-minimax-tree}中，root $A$ 为 MAX node，$B,C$ 为 MIN nodes，$D,E,F,G$ 为 MAX nodes。Terminal utilities 从左到右为
\[
  -1,\ 4,\ 2,\ 6,\ -3,\ -5,\ 0,\ 7.
\]
逐层计算各 internal node 的 Minimax value，并给出 $A$ 的选择。

\begin{checkedsolution}
从最底层 MAX nodes 开始：
\[
  \begin{aligned}
    D&=\max(-1,4)=4, & E&=\max(2,6)=6,\\
    F&=\max(-3,-5)=-3, & G&=\max(0,7)=7,\\
    B&=\min(4,6)=4, & C&=\min(-3,7)=-3,\\
    A&=\max(4,-3)=4.&&
  \end{aligned}
\]
MAX 应选择通向 $B$ 的 action。数值 $4$ 是面对 optimal opponent 时可保证的 utility，并非 opponent 犯错时的收益上限。
\end{checkedsolution}

\subsubsection*{同类检查}

若 MAX root 的三个 actions 分别进入 MIN nodes，terminal utilities 为
\[
  A_1:(3,1,8),\qquad A_2:(2,4,6),\qquad A_3:(14,5,2),
\]
则
\[
  V(A_1)=1,\qquad V(A_2)=2,\qquad V(A_3)=2.
\]
Root value 为 $2$，$A_2$ 与 $A_3$ 均为 optimal actions。

\textbf{检查点：}MIN 代表 opponent 的最优反制，不是随机选一个 child，因此不能对 children 求平均。

\clearpage
\exercisequestion{SC3000-L03-E02}{Game Show 的 Backward Induction}

\textbf{对应知识点：}
\begin{itemize}
  \item \relatedtopic{SC3000-L03-T04}{Return、Discount 与 Value Functions}
  \item \relatedtopic{SC3000-L03-T05}{Bellman Optimality Equation}
\end{itemize}

\subsubsection*{题目}

Game Show 共四题，答错后总奖金归零。依次答对的概率为 $0.9,0.75,0.5,0.1$；答完前三题后累计奖金分别为 $100,1100,11100$，第四题答对后为 $61100$。在每题之后可以 quit 并保留当前奖金。用 backward induction 决定策略并计算开始时的 expected value。

\begin{checkedsolution}
第四题继续作答的价值为
\[
  0.1(61100)+0.9(0)=6110<11100,
\]
所以答对第三题后应 quit。由后向前：
\[
  V(Q_3,\text{continue})=0.5(11100)=5550>1100,
\]
\[
  V(Q_2,\text{continue})=0.75(5550)=4162.5>100,
\]
\[
  V(Q_1,\text{continue})=0.9(4162.5)=3746.25.
\]
因此最优策略是继续回答前三题；若前三题均正确，则在第四题前退出。开始时的 expected value 为 $3746.25$。
\end{checkedsolution}

\textbf{检查点：}计算较早题目时，应使用后续 state 的 optimal value，而不是机械假设所有题都继续回答。

\clearpage
\exercisequestion{SC3000-L03-E03}{Terminal MDP 的 Expected Action Value}

\textbf{对应知识点：}
\begin{itemize}
  \item \relatedtopic{SC3000-L03-T03}{MDP、Markov Property 与 Policy}
  \item \relatedtopic{SC3000-L03-T05}{Bellman Optimality Equation}
\end{itemize}

\subsubsection*{题目}

在 state $s$ 中，action $a$ 以 $0.8$ 概率获得 reward $+1$ 并终止，以 $0.2$ 概率获得 reward $-1$ 并终止；action $b$ 确定获得 reward $+0.5$ 并终止。求两个 action values 并选择 optimal action。

\begin{checkedsolution}
所有 successor states 都是 terminal，因此 future value 为 $0$，discount factor 不影响本题：
\[
  Q(s,a)=0.8(1)+0.2(-1)=0.6,
  \qquad Q(s,b)=1(0.5)=0.5.
\]
故
\[
  V^*(s)=\max(0.6,0.5)=0.6,
  \qquad \pi^*(s)=a.
\]
\end{checkedsolution}

\textbf{检查点：}Agent 对可选 actions 取 $\max$；同一 action 的随机 outcomes 必须按 transition probabilities 求 expectation。
